\chapter{Trust Establishment}
\label{chap:trust-establishment}

\begin{abstract}
Parties cannot remove trust by naming a cryptographic primitive.  They can
instead make each trust assumption visible and ask which property depends on
it.  This chapter separates local entropy generation, public coordination,
unauthenticated key exchange, identity binding, authenticated key
establishment, and secret-key agreement from correlated observations.  Its
central boundary is simple: local randomness may create a private secret, but
it does not identify a remote participant.  The chapter presents a taxonomy
and research agenda, not a new key-establishment protocol or security theorem.
\end{abstract}

\section{Scope and Contribution Status}

Within this book, \emph{trust establishment} means making explicit which setup
resource binds a session, credential, or key to an intended peer, and which
properties that binding supplies.  It does not mean creating moral trust.  It
does not prove that a participant is honest, that an institution assigned the
right identity, or that messages will be delivered.

This chapter is a literature-backed foundation.  It makes no novelty claim,
selects no universal identity system, and implements no protocol.  Its task is
to prevent later chapters from confusing distinct resources merely because
they all appear near the beginning of a cryptographic conversation.

\section{Author's Intuition}

The original draft begins with an important instinct:

\begin{quote}
Each participant should isolate and safeguard an experiment from which local
uncertainty is measured.  The participant should not have to surrender every
private advantage merely to cooperate.
\end{quote}

That instinct survives here as a principle of participant autonomy.  A party
should be able to control its local source, understand its failure modes, and
retain only the secrets the joint purpose actually needs.  The draft's phrase
``secrets are tactical disadvantages'' is therefore preserved as a
secret-minimization principle: long-lived secrets are capabilities, but they
are also liabilities that can be stolen, coerced, copied, or misused.

The earlier text also distinguishes an adversary, who may protect its own
interests while sharing the joint purpose, from an enemy whose objective
includes defeating that purpose.  This chapter preserves the distinction as an
author's modeling question, not as a cryptographic property.  A later protocol
must express it through explicit behavior or preference assumptions.

It must not be read as the technical claim that secret-key disclosure causes
no harm.  Disclosure can enable impersonation, forgery, decryption, or
retrospective compromise, depending on the protocol.

An earlier version of this chapter also asks whether extensive challenge steps
can defeat an active interposer, calls the problem ``a grail in cryptography,''
and separates active interception from denial of service.  Those questions
remain valuable.  A challenge, however, can demonstrate possession only
relative to a credential, secret, or statement that was already bound to the
intended peer.  Challenges do not create that binding from nothing.

\section{The Cooperation Problem}

Let Alice and Bob intend to communicate.  Let Mallory control the network, and
let $\Sigma$ denote whatever setup resource the parties already possess.  The
resource $\Sigma$ might be empty, a pre-shared secret, an authenticated public
key, a certificate policy, a remembered first key, a password, or evidence
obtained in a physical meeting.  Let $T$ be the public transcript.  Depending
on the task, Alice and Bob may output local keys $K_A$ and $K_B$, a public
action, or an explicit abort.

The phrase ``channel security'' hides several different questions:

\begin{itemize}
    \item \textbf{Confidentiality:} who can read the messages?
    \item \textbf{Origin authentication:} what justifies attributing a message
    to a particular peer or credential?
    \item \textbf{Integrity:} can modification be detected?
    \item \textbf{Freshness:} can an old valid message be replayed as new?
    \item \textbf{Ordering and agreement:} do the parties accept the same
    session context and transcript?
    \item \textbf{Delivery:} can an adversary delay or suppress every message?
\end{itemize}

A chapter may assume some of these properties and construct others.  It may
not silently treat them as one indivisible guarantee.

\section{Six Different Tasks}

The comparison below separates six tasks that the earlier draft placed too
close together.

\begin{center}
\centering
\small
\textbf{Six tasks and their resources}\par\medskip
\begin{tabular}{>{\raggedright\arraybackslash}p{0.20\textwidth}
                >{\raggedright\arraybackslash}p{0.24\textwidth}
                >{\raggedright\arraybackslash}p{0.22\textwidth}
                >{\raggedright\arraybackslash}p{0.22\textwidth}}
\textbf{Task} & \textbf{Initial resource} & \textbf{Output or goal} &
\textbf{What is still missing} \\
\hline
Local secret generation & Characterized local source and protected endpoint &
Unpredictable local bits & No remote peer or shared value \\
\hline
Public coordination & Common public state, algorithm, or beacon & Same public
action or random value & No secrecy and no identity by itself \\
\hline
Unauthenticated key exchange & Computational assumptions and fresh local
randomness & Matching key against a passive observer & No protection from
active key substitution \\
\hline
Identity or credential binding & A setup rule in $\Sigma$ & Association between
a name or role and a credential & No session key or delivery by itself \\
\hline
Authenticated key establishment & A binding resource plus a key-establishment
protocol & A fresh, peer-bound shared key or abort & Availability and endpoint
integrity remain external \\
\hline
Correlated-source agreement & Jointly distributed observations $X,Y,Z$ and
authenticated public discussion & Information-theoretic shared key under an
exact source model & No conclusion from Alice--Bob correlation alone \\
\end{tabular}
\end{center}

The table is not a menu of interchangeable implementations.  It is a guard
against circular reasoning.  For example, assuming an authenticated public
channel in order to study correlated-source agreement is legitimate, but that
study does not then explain where the authentication came from.

\section{Local Entropy Is Not Remote Identity}

A defensible local randomness pipeline contains several separately reviewable
stages:

\begin{enumerate}
    \item describe the physical or non-physical noise source and how it may be
    observed or influenced;
    \item digitize raw observations without hiding finite-precision or sampling
    effects;
    \item estimate a conservative lower bound on min-entropy under a stated
    model;
    \item condition or extract only under the preconditions of a named
    construction;
    \item derive purpose-specific keys with explicit context and separation;
    \item health-test the source, protect the resulting state, rotate secrets,
    and define failure behavior.
\end{enumerate}

NIST SP 800-90B treats a noise source, optional deterministic conditioning,
and health testing as distinct components.  It requires documentation and
entropy assessment and states that deterministic conditioning has at most the
entropy of its input\cite{nist80090b}.  Conditioning may reduce bias or
concentrate available entropy; it cannot manufacture entropy absent from the
source.  A later key-derivation step likewise does not repair an
uncharacterized input merely by giving it a cryptographic name.

Covariance, variance, visible noise, and an attractive histogram can help
diagnose a source, but none is automatically a min-entropy bound.  Physical
shielding may reduce some forms of observation or influence, yet it proves
neither source quality nor endpoint integrity.  If Mallory can predict the
source, bias the measurement, or read the conditioned key from memory, the
local pipeline has failed even if every deterministic transformation executes
correctly.

Most importantly, Alice can execute a flawless local pipeline and still have
no evidence that a received message came from Bob.  Local unpredictability and
remote identity answer different questions.

\subsection{Open Research Note: The Original Entropy Model}

The original sketch proposed raw observations $\Pi_n$, a participant-specific
model $\Lambda$, covariance eigenvalues, a quantizer into $\mathbb{Z}_p$, and a
final conditioning step.  It also wanted participants to model local noise
according to their own preferences.  The sketch was explicitly unfinished.

It is preserved as a research question rather than an active protocol.  Before
it could derive cryptographic key material, it would need at least:

\begin{itemize}
    \item a defined stochastic and adversarial source model;
    \item an explanation of what $\Lambda$ maps and what its eigenvalues mean;
    \item a conservative min-entropy estimate with confidence and restart
    behavior;
    \item a finite-precision quantizer and an analysis of bias after mapping
    into $\mathbb{Z}_p$;
    \item exact preconditions for the conditioner or extractor; and
    \item health tests and explicit behavior when the model is falsified.
\end{itemize}

Two questions are especially promising.  Can covariance structure assist
source diagnostics without being mistaken for an entropy proof?  Can parties
choose autonomous local sources and models while still reporting an objective,
comparable lower bound required by a shared protocol?  Until those questions
are answered, the construction remains an open research note.

\section{From Key Exchange to an Intended Peer}

An unauthenticated computational key exchange can be analyzed against a
passive observer under the assumptions of its selected construction.  That
statement does not extend to an active network adversary who may replace
messages.

The elementary interposition trace is:

\begin{enumerate}
    \item Alice sends unauthenticated key-exchange material intended for Bob.
    \item Mallory replaces it and sends Mallory's material to Bob.
    \item Bob responds; Mallory replaces Bob's material before Alice sees it.
    \item Alice establishes one key with Mallory, while Bob establishes a
    different key with Mallory.
    \item Mallory may relay, modify, or suppress subsequent traffic.
\end{enumerate}

Mallory, as an endpoint of both substituted exchanges, knows both resulting
keys.  A separate passive observer may still find the keys unpredictable, but
that passive-observer property binds neither key to the intended peer.  This is
why local randomness is necessary for many protocols but insufficient for peer
authentication.

TLS 1.3 provides one current standards example of the separation: its
handshake has distinct key-exchange and authentication work, can use
certificate-based credentials or a pre-shared key, and provides key
confirmation and handshake integrity.  Its unauthenticated modes still require
an external binding such as out-of-band validation or trust on first use to
resist active interposition\cite{rfc9846tls13}.  TLS is an example here, not a
protocol proposed or implemented by this book.

\subsection{What a Challenge Can Establish}

A challenge/response exchange can demonstrate that the responder possesses a
secret key, password-derived capability, witness, or hardware credential.  It
does not answer why that capability should be called ``Bob.''  That answer
comes from $\Sigma$: a prior ceremony, credential issuer, pin, remembered key,
or other binding rule.

A complete session also needs context.  The proof or authenticator should bind
the protocol and version, parties or roles, session identifier, transcript,
and freshness material.  Otherwise a valid response may be replayed or moved
into a different conversation.  Key confirmation then establishes that a peer
possessing the relevant secret derived the same session key; it still does not
guarantee that messages will arrive or that the endpoint is uncompromised.

\section{Setup Choices Are Different Resources}

No single setup choice is universally correct.  The comparison below records
the most important tradeoffs without choosing an implementation.

\noindent\textbf{Pre-shared high-entropy secret.}
May bind peers that received the same secret.  Principal risks are secure
distribution, copying, group ambiguity, and compromise.

\medskip
\noindent\textbf{Authenticated public key, certificate, or pin.}
May bind a name or role under a verification policy.  Principal risks are
issuer or policy error, revocation failure, key theft, and pinning mistakes.

\medskip
\noindent\textbf{Trust on first use (TOFU).}
May bind later sessions to the first observed key.  The first contact can be
captured, and legitimate key change is difficult.

\medskip
\noindent\textbf{Password-authenticated exchange.}
May bind a peer that knows the agreed password under the chosen protocol.
Principal risks include weak passwords, enrollment, recovery, coercion, and
protocol-specific assumptions.

\medskip
\noindent\textbf{Human comparison or ceremony.}
A fingerprint comparison or physical ceremony may bind a short value or
credential to people present.  Principal risks are human error, accessibility,
coercion, and ceremony substitution.

A credential therefore authenticates only relative to its issuance, storage,
verification, and revocation policy.  A valid certificate can show that a
configured authority made a statement.  It cannot show that the authority's
statement is just, that a human remains uncoerced, or that the credential's
private key was never copied.

\section{Correlated Observations and Public Discussion}

Information-theoretic secret-key agreement from observations requires a model
that includes the eavesdropper.  For the cited rate statement, fix a
finite-alphabet joint distribution $P_{XYZ}$ and give the parties $N$
independent repetitions $X^N$, $Y^N$, and $Z^N$.  Let $T_N$ contain every
message exchanged during the public-discussion protocol.  The asymptotic rate
counts secret-key bits per repetition as $N$ tends to infinity, under vanishing
disagreement probability and vanishing normalized information leakage to Eve
given $(Z^N,T_N)$.  Secrecy therefore cannot be inferred merely from the
correlation between $X$ and $Y$.

Maurer's repeated-source model assumes that active modification or message
injection on the public discussion channel is detected.  In that model, the
secret-key rate satisfies\cite{maurer1993secretkey}

\[
    S(X;Y\mathbin\Vert Z)
    \leq \min\!\left\{I(X;Y),\,I(X;Y\mid Z)\right\}.
\]

Thus, if $X$ and $Y$ are independent, the information-theoretic secret-key
rate is zero in this model.  This is a conditional result, not a general claim
that computational key agreement is impossible.  Public-key assumptions can
change what computationally bounded parties can achieve.

Positive mutual information between Alice and Bob is also insufficient.  Let
$H(X)>0$ and $X=Y=Z$.  Then $I(X;Y)=H(X)>0$, but Eve already knows the entire
common observation.  Likewise, error-correction or reconciliation messages in
$T_N$ may reveal
information that must be accounted for before privacy amplification.  A rate
slogan that omits $Z$, $T_N$, agreement error, and the exact source distribution
does not establish a secret key.

This distinction repairs the earlier statement that any correlated common
phenomenon can be ``squeezed'' into secrecy.  Correlation may be useful, but its
value is always relative to Eve's observation and the public leakage required
to reconcile Alice's and Bob's views.

\section{Public Coordination Without Secrecy}

Public common randomness is valuable precisely because many observers can
verify and reuse it.  NIST describes randomness beacons as a public utility for
auditable randomized procedures and explicitly warns that beacon values must
not be used as secret cryptographic keys\cite{nist-randomness-beacons}.  Alice,
Bob, and Eve can all agree on the same beacon value; agreement is not secrecy.

The same distinction applies to deterministic optimization.  If Alice and Bob
have identical public input, the same specified algorithm and version, and
deterministic tie-breaking, they obtain the same public output.  For example,
two routers using the same topology snapshot, weights, Dijkstra implementation,
and tie-breaking rule can select the same route.  An observer with the same
information can compute it too.

The original note that ``a market is a shared channel'' is preserved as an
open empirical hypothesis, not as a secret-key claim.  A serious formulation
must define what Alice observes as $X$, what Bob observes as $Y$, what Eve sees
as $Z$, what public interaction becomes $T$, and which observations are not
already public.  Market correlation alone does not answer those questions.

\section{Attacks and Counterexamples}

The following cases should be tested against any future construction:

\begin{description}
    \item[Active interposition.] Mallory establishes separate sessions with
    Alice and Bob when exchanged keys lack an authentic binding.
    \item[Credential substitution.] Mallory supplies Mallory's public key under
    Bob's name when no trusted association exists.
    \item[First-contact capture.] A TOFU system faithfully remembers the
    attacker's key because the first session was already intercepted.
    \item[Replay.] A valid old proof is accepted in a new session because the
    authenticator omitted freshness or context.
    \item[Source manipulation.] Mallory biases or predicts the local source;
    deterministic conditioning cannot reconstruct missing entropy.
    \item[Quantizer bias.] A nonuniform or finite-precision source maps
    disproportionately into some residues.
    \item[Endpoint compromise.] Mallory reads the final key directly from
    memory, bypassing the mathematical pipeline.
    \item[Key exposure.] A stolen pre-shared or signing key enables
    impersonation and may compromise protected traffic according to the
    protocol.
    \item[Public correlation.] Let $H(X)>0$ and $X=Y=Z$.  Alice and Bob have
    positive mutual information and perfect agreement, while Eve has the same
    information.
    \item[Reconciliation leakage.] Public correction data removes conditional
    entropy that an informal calculation counted as private.
    \item[Public beacon.] All parties agree, including Eve; no shared secret is
    created.
    \item[State divergence.] Different router state or tie-breaking produces
    different actions.
    \item[Denial of service.] Mallory drops every message.  Without an explicit
    timing or failure-detector assumption, suppression may be indistinguishable
    from delay.  Alice may observe only local non-completion or a timeout;
    cryptography neither identifies the cause nor forces delivery.
\end{description}

\section{What Cryptography Does Not Solve}

Cryptography can bind statements to keys, keys to transcripts, and transcripts
to an explicitly configured setup.  It cannot decide whether the configured
identity policy is fair, whether a certificate issuer spoke truthfully,
whether a participant is coerced, or whether cooperation is wise.  It cannot
make a failed entropy source unpredictable, protect a key already read from a
compromised endpoint, or compel an adversary to deliver a message.

This chapter therefore establishes no interpersonal or moral trust, truthful
input, availability, fairness, non-repudiation, coercion resistance,
composability, post-quantum guarantee, or general information-theoretic
guarantee.  It establishes no secret key at all.  Its contribution is the
boundary map that later protocol chapters must fill with an explicit setup,
adversary, construction, and evidence.

\section{Limitations and Open Problems}

The taxonomy is deliberately prior to construction.  Important open problems
remain:

\begin{enumerate}
    \item Choose one target setting for this book and state how human or
    institutional identity is bound to protocol credentials.
    \item Determine whether the original covariance model can support source
    diagnostics and a defensible min-entropy bound.
    \item Specify one physical or non-physical source, its adversarial influence
    model, its quantizer, health tests, and its failure behavior.
    \item Decide when TOFU, a pre-shared secret, a password ceremony, a
    certificate policy, or a physical meeting is an acceptable setup.
    \item Define a market-observation experiment with explicit $X,Y,Z,T$ and
    measure whether any conditional secrecy remains after public discussion.
    \item Separate detection and accountability from recovery when an active
    adversary blocks the channel.
\end{enumerate}

\section{Exercises}

\begin{enumerate}
    \item For a proposed messaging protocol, list separately its assumptions
    about confidentiality, authentication, integrity, freshness, ordering, and
    delivery.
    \item Draw the two-session man-in-the-middle trace and identify the first
    message that would need an authentic binding to prevent substitution.
    \item Explain why a challenge signed by an unknown key proves possession
    of that key but not the signer's intended human identity.
    \item Give a concrete nonconstant example with $H(X)>0$ and $X=Y=Z$,
    then compute why Alice--Bob correlation alone cannot imply secrecy from
    Eve.
    \item Compare a pre-shared secret, TOFU, and a physical fingerprint
    ceremony.  State one failure unique to each.
    \item Construct two public shortest-path inputs that differ by one edge or
    one tie-breaking rule and cause two deterministic routers to disagree.
    \item Propose a measurable hypothesis for the market-observation idea,
    including Eve's view and every public reconciliation message.
\end{enumerate}
